% !PRINTED 331
\weberpage{331}

\gesperrt{Unter der Norm} $N(\mathfrak{a})$ des Ideals $\mathfrak{a}$ versteht man den
\gesperrt{absoluten Wert} der Determinante $A$ der Substitution $(A)$:
\[
A = \Sigma \pm a_{1,1}\,a_{2,2}\ldots a_{n,n},
\]
und nach (3) ist die Norm unabhängig von der Wahl der Basis.

1. \gesperrt{Wir wollen die Basis $(\alpha_\nu)$ positiv nennen, wenn
die Determinante $A$ positiv, also der Norm von $\mathfrak{a}$ gleich ist.
Ist $(\alpha_\nu)$ positiv, so ist $(\alpha'_\nu)$ immer dann
positiv, wenn die Substitutionsdeterminante $L = +1$ ist.}

Man kann aus jeder Basis durch eine Substitution mit der
Determinante $-1$, also z.\ B. durch Vertauschung zweier $\alpha_\nu$, eine
positive Basis ableiten, und wir werden in der Folge meist nur
positive Basen verwenden.

Die in Bd. II, \S\ 163, (3) bestimmte Basis \gesperrt{ist positiv}.

Bedeutet nun $t_1, t_2, \ldots, t_n$ ein System unabhängiger Vari\-ablen und
\weq{4}{\lambda = \sum_\nu \alpha_\nu t_\nu}
% Source scan confirms nu as the lower summation index in (4).
eine Basisform mit positiver Basis, so ergibt sich (Bd. II, \S\ 164)
\weq{5}{N(\lambda) = N(\mathfrak{a})\,T,}
worin
\weq{6}{T = \Sigma \pm t_{1,1}\,t_{2,2}\ldots t_{n,n}}
eine \gesperrt{primitive ganzzahlige homogene Form} $n$ten Grades der
Variablen $t_\nu$ ist. Die $t_{r,s}$ sind \gesperrt{lineare Funktionen} der $t_\nu$.

Wendet man auf (4) die Substitution (2) an, so ergibt sich
\weq{7}{\lambda = \lambda' = \sum \alpha'_\nu t'_\nu,}
wenn
\weq{8}{(t_1, t_2, \ldots, t_n) = L'(t'_1, t'_2, \ldots, t'_n),}
worin $L'$ die \gesperrt{transponierte Substitution} von $L$ ist, und wenn
$(\alpha'_\nu)$ gleichfalls eine positive Basis ist, so hat $L'$ die
Determinante $+1$ und die Formel (5) ergibt
\weq{9}{N(\lambda) = N(\mathfrak{a})\,T',}
worin $T'$ eine homogene Funktion $n$ten Grades der Variablen $t'$
ist, die durch die Substitution (8) aus $T$ hervorgeht.

Zwei Formen, die durch ganzzahlige lineare Substitution mit
der Determinante $+1$ auseinander hervorgehen, heißen \gesperrt{äquivalent}. Bisweilen werden die Formen auch dann äquivalent
genannt, wenn die Substitutionsdeterminante $-1$ ist, dann aber
mit dem Zusatze \glqq \gesperrt{uneigentlich}\grqq{}. Es ist dann durch (8) bewiesen:
%CONT>>
